\usepackage{array}
\usepackage{hyperref}
\usepackage{graphicx} % Allows including images
\usepackage{booktabs} % Allows the use of \toprule, \midrule and \bottomrule in tables
\usepackage{graphicx}
\usepackage{amsmath}
\usepackage{amssymb}
\usepackage{physics}
\usepackage{braket}
\usepackage{mathpartir}
\usepackage{pdfpages}

\usepackage{glossaries}
% \makeglossaries
\newacronym{ra}{RA}{Register Assignment}
\newacronym{jair}{JAIR}{Just Another Intermediate Representation}
\newacronym{ssa}{SSA}{Static Single Assignment}
\newacronym{peo}{PEO}{Perfect Elimination Order}
\newacronym{cfg}{CFG}{Control Flow Graph}
\newacronym{bnf}{BNF}{Backus-Naur Form}
\newacronym{nasm}{NASM}{Netwise Assembler}
\newacronym{jit}{JIT}{Just-In-Time}
\newcommand{\nt}[1]{\langle\text{#1}\rangle}
\newcommand{\BareCoqSymbol}{\includegraphics[height=0.9em]{coq.pdf}}
\newcommand{\CoqSymbol}{\raisebox{-.2ex}{\BareCoqSymbol\,}}
\newcommand{\Coqed}{{\hfill\CoqSymbol}}
\newcommand{\chordalpeosymbol}{\mathrm{Chordal}_{\gls{peo}}}
\newcommand{\chordalindsymbol}{\mathrm{Chordal}_{ind}}
\newcommand{\chordalpeo}[1]{\vdash {#1} : \chordalpeosymbol}
\newcommand{\chordalind}[1]{\vdash {#1} : \chordalindsymbol}
\newcommand{\simplicialsymbol}{\mathrm{Simplicial}}
\newcommand{\simplicial}[2]{{#2} \vdash {#1} : \simplicialsymbol}
\newcommand{\addvertex}[3]{#1 \oplus_{#2} #3}
\newcommand{\colorline}[2]{%
\begingroup
\setlength{\fboxsep}{0.5pt}%
\colorbox{#1}{#2}%
\endgroup
}
\newcommand{\backupbegin}{
   \newcounter{framenumberbeforeappendix}
   \setcounter{framenumberbeforeappendix}{\value{framenumber}}
}
\newcommand{\backupend}{
   \addtocounter{framenumberbeforeappendix}{-\value{framenumber}}
   \addtocounter{framenumber}{\value{framenumberbeforeappendix}}
}

\usepackage{tikz}
\usetikzlibrary{arrows.meta, positioning, calc}
\usetikzlibrary{shadows}
\usetikzlibrary{shadows.blur}

\newcommand{\contr}[1]{build/contr#1.pdf}
\newcommand{\trycontr}[1]{
  \IfFileExists{\contr #1}{
    \includegraphics[page=1,width=0.4\textwidth]{\contr #1}
  }{
    {\color{red} [Slide thumbnail will appear after next compile]}
  }
}

\usepackage{listings}
\lstset{
    basicstyle=\ttfamily,  % Same font as alltt
    columns=fullflexible,  % Match spacing of alltt
    keepspaces=true,       % Preserve spaces like alltt
    xleftmargin=2em,
}

\definecolor{vara}{RGB}{255,102,102}  % Red for variable a
\definecolor{varb}{RGB}{102,178,255}  % Blue for variable b
\definecolor{varc}{RGB}{102,255,102}  % Green for variable c
\definecolor{vard}{RGB}{255,178,102}  % Orange for variable d

% Colorblind-proof colors
\definecolor{StrongOrange}{RGB}{230,159,0}
\newcommand{\cbnodeA}[0]{RoyalBlue}
\newcommand{\cbnodeB}[0]{RedOrange}
\newcommand{\cbnodeC}[0]{StrongOrange}

\definecolor{myblue}{RGB}{0,0,180}
\definecolor{mygreen}{RGB}{0,120,0}
\definecolor{myred}{RGB}{180,0,0}
\definecolor{mypurple}{RGB}{120,0,120}

\lstdefinestyle{C}{
    language=C,
    mathescape=true,
    tabsize=4,
    keywordstyle=\color{myblue},              % Keywords in blue
    commentstyle=\color{mygreen!50!black},    % Comments in green
    stringstyle=\color{orange},             % Strings in orange
    showstringspaces=false,
    numbers=left,
    numberstyle=\tiny\color{gray},
    stepnumber=1,
    % frame=single,
    breaklines=true
}

\lstdefinelanguage{Rocq}{
    morekeywords=[1]{Definition, Inductive, CoInductive, Class, Instance, Fixpoint, CoFixpoint, Function, Proof, Qed, Module, Fail},
    morekeywords=[2]{forall, exists, match, with, end, if, then, else, let, in, fix, fun, intros, unfold, destruct, lia},
    morekeywords=[3]{Prop, Type, Set, True, False, nat, bool, option, list},
    sensitive=true,
    morecomment=[l]{(*},
    morecomment=[s]{(*}{*)},
    morestring=[b]",
}

\lstdefinestyle{Rocq}{
    language=Rocq,
    mathescape=true,
    keywordstyle=[1]\color{myblue},
    keywordstyle=[2]\color{mypurple},
    keywordstyle=[3]\color{myred},
    tabsize=2,
    commentstyle=\color{green!50!black}\itshape,
    stringstyle=\color{orange},
    showstringspaces=false,
    numbers=left,
    numberstyle=\tiny\color{gray},
    stepnumber=1,
    % frame=single,
    breaklines=true
}

% OCaml language with keyword groups
\lstdefinelanguage{OCaml}{
  sensitive=true,
  morecomment=[s]{(*}{*)},
  morestring=[b]",
  morekeywords=[1]{
    let, rec, in, and, val, type, module, struct, sig, class, object, method,
    inherit, virtual, functor, include, open, external, exception
  },
  morekeywords=[2]{
    if, then, else, match, with, try, when, for, to, downto, while, do, done,
    begin, end, fun, function
  },
  morekeywords=[3]{true, false},
  morekeywords=[4]{of, mutable, private, new, lazy, as, assert}
}

\lstdefinestyle{OCaml}{
    language=OCaml,
    mathescape=true,
    keywordstyle=[1]\color{myblue}\bfseries,    % Declarations
    keywordstyle=[2]\color{mypurple}\bfseries,  % Control flow
    % keywordstyle=[3]\color{orange}\bfseries,    % Constants
    keywordstyle=[4]\color{myred}\bfseries,      % Misc keywords
    tabsize=2,
    basicstyle=\ttfamily\small,
    commentstyle=\color{green!50!black}\itshape,
    stringstyle=\color{orange},
    showstringspaces=false,
    numbers=left,
    numberstyle=\tiny\color{gray},
    stepnumber=1,
    % frame=single,
    breaklines=true,
}

\lstdefinelanguage{NASM}{
  morekeywords={
    mov, add, sub, div, mul, imul, idiv, xor, and, or, inc, dec, jmp, cmp, je, jne, jl, jle, jg, jge, call, ret, push, pop, syscall
  },
  sensitive=true,
  morecomment=[n]{;}{}, % NASM uses ';' for comments
  morestring=[b]",      % Strings in double quotes
  morestring=[b]',      % Strings in single quotes
}

\lstdefinestyle{NASM}{
  language=NASM,
  keywordstyle=\color{myblue}\bfseries,
  numbers=left,
  numberstyle=\tiny\color{gray},
  stepnumber=1,
  % frame=single,
  basicstyle=\ttfamily,
  breaklines=true
}

